\subsection{長さ二・一・一の表示と正規元による縮小}
\label{subsec:as4-ds211-exclusions}
引き続きas4の固定commit\cite{AS4lowpinned}から具体的表示だけを参照する。
本節の最初の二族は
{\small\texttt{CY\_ZERO\_COMMUTATORS\_DOUBLE\_SIMPLE\_SECOND\_COMPLETE.tex}}
の(DT1)である。

\begin{theorem}[二パラメータOre表示の全排除]
\label{thm:ds211-A-all-excluded}
$l\ne0,\pm1/4$、$r\in k$ とし
\[
 q=\frac{4l+1}{4l-1},\quad \beta=-\frac1{4l},\quad
 u=-\frac{16lr+1}{2l(16l^2-1)},\quad v=\frac{8(l+r)}{16l^2-1}
\]
と置く。表示 $\mathcal A_{l,r}$ は量子多項式環
\[
 S=k\langle b,c,d\mid cb=qbc,\ db=q^{-1}bd,\ dc=q^2cd\rangle
\]
上のOre拡大 $S[a;\sigma,\delta]$ であり
\[
 \sigma(b,c,d)=(b,q^{-1}c,qd),\qquad
 \delta(b,c,d)=(\beta b^2,u bc,v bd)
\]
となる。全パラメータでH3包絡にはならない。
中心三次が非零となるのは
\[
 48l^2+1=0,\quad r=3l\quad\text{または}\quad r=-l
\]
だけであり、その中心三次は順に $kc^3$、$kd^3$ である。
\end{theorem}
\begin{proof}
\texttt{verify\_cubic\_as4\_ds211.py}は元の六関係と上の六Ore関係の一致、
順序 $b<c<d<a$ の四つのPBW重なりを全係数で検証する。
従って $S$ 上の多項式正常形があり、重み $w(a)=2,w(b,c,d)=1$ の先頭環が量子四平面なので整域である。

$b$ のnormalizerは
\[
 a\longmapsto a-\beta b,\quad b\longmapsto b,\quad
 c\longmapsto q^{-1}c,\quad d\longmapsto qd
\]
である。その半単純部分と単冪部分は可換し、単冪部分は $a\mapsto a-\beta b$、他の三生成元を固定する。
これらは直接Ore関係を保つ自己同型である。
中心三次 $F$ はnormalizerで固定されるので、その有限次元作用のJordan分解から
単冪部分でも固定される。
$F=\sum_{j=0}^m f_j a^j$ と正常形で書くと、$m>0$ の場合
単冪変換との差の最高 $a$ 次数係数は $-m\beta f_m b$ となり非零である。
従って全中心三次は $S_3$ に入る。

$S$ の単項式 $b^ic^jd^k$ が中心である条件は
$q^{j-k}=q^{i-2k}=1$ である。
$i+j+k=3$ の十組について、二指数の最大公約数は純三乗では三、それ以外では一である。
$q\ne1,-1$ なので、$q^3\ne1$ では中心三次は零、$q^3=1$ では
$S$ の中心三次は $b^3,c^3,d^3$ の三本に限る。
このとき
\[
 [a,b^3]=3\beta b^4,\quad
 [a,c^3]=3u bc^3,\quad [a,d^3]=3v bd^3.
\]
これらも検証器で三乗根の係数環上で厳密確認した。
従って $b^3$ の係数は零、$c^3,d^3$ はそれぞれ $u=0,v=0$ のときだけ残る。
$u=v=0$ は $16l^2=1$ を要求するので起こらない。
$q^3=1$ は $48l^2+1=0$ に等しく、その上で $u=0$ は $r=3l$、$v=0$ は $r=-l$ となる。

H3に必要な中山作用も、AS4実現を仮定したときに独立に恒等と判定できる。
対称射影を $U=(b^2,\langle a,c\rangle,\langle a,d\rangle,
\langle b,c\rangle,\langle b,d\rangle,\langle c,d\rangle)$ とする正規化関係行列 $R$ と
\[
 A_1=\tfrac12\operatorname{Hess}(-z_1z_4+z_2z_3-z_0z_5+z_3z_4)
\]
を使い、$RA_1R^{\mathsf T}$ の非零性と全256係数の反巡回性を確認した。
前節と同じ最高双対汎関数の議論から $\mu=I$ である。
従ってH3候補は上の中心三次に限り、残る純三乗も非整域商を与える。
\end{proof}

\begin{theorem}[例外的一パラメータ表示の全排除]
\label{thm:ds211-B-all-excluded}
$i^2=-1$ とし、$\mathcal B_r$ を次の六関係で定める。
\[
\begin{split}
0&=b^2+i[a,b]-2[c,d],\\
0&=(1+i)ac+(1-i)ca-(4r-i)[b,c],\\
0&=(1-i)ad+(1+i)da+(4r+i)[b,d],\\
0&=(1+i)bc+(1-i)cb,\\
0&=(1-i)bd+(1+i)db,\qquad0=cd+dc.
\end{split}
\]
全 $r\in k$ でこの表示はH3包絡にならない。
\end{theorem}
\begin{proof}
前証明と同じ $A_1$ と正規化関係行列を用い、非零最高tensorの反巡回性から
AS4実現における $\mu=I$ を独立に確認する。
$c,d$ のnormalizerは
\[
\begin{array}{c|rrrr}
 &a&b&c&d\\\hline
\nu_c&-ia+(1+4ir)b&-ib&c&-d\\
\nu_d&ia+(4ir-1)b&ib&-c&d
\end{array}
\]
であり、全八式を元の二次関係の線型結合として確認した。
従って $N=c^3d$ のnormalizerは
\[
 a\longmapsto-a+(8r-4i)b,\quad b\longmapsto-b,\quad
 c\longmapsto-c,\quad d\longmapsto-d.
\]
半単純部分は $-I$ なので、H3包絡を仮定した整域では
補題\ref{lem:normalizer-minus-unipotent-odd-centre}が中心三次を消す。
これで全 $r$ を排除する。
\end{proof}

\begin{theorem}[残る二族と固定表示の三次核消滅]
\label{thm:ds211-TU-fixed-excluded}
同commitの
{\small\texttt{CY\_\allowbreak{}ZERO\_\allowbreak{}COMMUTATORS\_\allowbreak{}DOUBLE\_\allowbreak{}SIMPLE\_\allowbreak{}NONZERO\_\allowbreak{}COMPLETE.tex}}
の(DS3)にある $\mathcal T_h$ ($h\ne0$)、$\mathcal U_u$ ($u\ne0$)、
$\mathcal L_{i/4},\mathcal L_{-i/4},\mathcal P$ は全てH3包絡にならない。
また同表の $\mathcal V_0$ も中心三次を持たない。
\end{theorem}
\begin{proof}
\texttt{verify\_cubic\_as4\_ds211\_remaining.py}は(DS1)--(DS3)から全六関係を再構成する。
最高tensorの全256反巡回恒等式を確認し、AS4実現で $\mu=I$ となることを独立に得る。
$\mathcal T_h,\mathcal U_u$ の中心三次行列から選んだ最大minorはそれぞれ
\[
 \frac{(8h+1)(16h+1)}{2^{23}h^6},\qquad
 -2^{29}u^6(4u-1)(48u^2+8u+1)
\]
である。次数三・四の商座標の分母は前者では $h$ の冪と有理定数だけ、後者では有理定数だけである。
商座標は全関係を消し、選んだ自由列上では恒等である。
従って任意のAS4実現では、次数三・四の次元 $20,35$ を使えばこれらは全許容値で実際の商基底となる。
一般点のrow reductionの結果だけから、特殊fiberの階数不変性を推論してはいない。

残る値は $h=-1/8,-1/16$、$u=1/4$ および $48u^2+8u+1=0$ だけである。
最後の二値は $u=(-1\pm\sqrt{-2})/12$ と書ける。
有理数または $\mathbb Q(\sqrt{-2})$ 上で、各値の元の全関係商を改めて計算し、
全てで中心三次行列の階数が $20$ となる非零minorを得た。
数体の二つの埋込みが両根を覆うので、有限標本からの外挿ではない。
これで $\mathcal T_h,\mathcal U_u$ の全許容値を覆う。

固定表示 $\mathcal L_{i/4},\mathcal P$ は $\mathbb Q(i)$ 上、
$\mathcal V_0$ は $\mathbb Q$ 上で同じ階数 $20$ を得る。
$\mathcal L_{-i/4}$ は係数体の共役により同じ結論となる。
全六関係、商座標の分母、minorの行番号と厳密な行列式をコードと証明書に保持した。
従って全ての表示で必要な中心三次が存在しない。
\end{proof}

\begin{theorem}[最後の一パラメータ表示 $\mathcal V_h$ の全排除]
\label{thm:ds211-V-all-excluded}
(DS3)の $\mathcal V_h$ ($h\ne-1/16$) も全てH3包絡にならない。
その中心三次は $h\ne1/48$ では零、$h=1/48$ では四次元であり、
後者の全非零切断も一次式と二次式の積に分解する。
従って(DS3),(DT1)に列挙された全表示がH3包絡から排除される。
この一覧が全AS4の当該射影型を覆うという定理そのものは、ここでは入力としていない。
\end{theorem}
\begin{proof}
$h=0$ は前定理で処理した。
$h\ne0$ では $r^2=-h$ を取り、$r\ne0,\pm1/4$ とする。
\[
 p=\frac{4r-1}{4r+1},\quad
 \alpha=\frac{p^2}{2(p+1)^2},\quad
 \beta=\frac1{2(p+1)^2},\quad \gamma=\frac{p}{(p+1)^2},
\]
\[
 a'=a+\gamma b,\quad x=c-rd,\quad y=c+rd,\quad
 \vartheta=p-p^{-1},\quad
 t_0=\frac{p-1}{2(p+1)},\quad
 s_0=\frac{(p-p^{-1})(1+p^2)}{2(p+1)^2}
\]
と置く。この有限平方根被覆は $h\ne0,-1/16$ を漏れなく覆い、$p\ne0,\pm1$ である。
新しい六関係は
\[
\begin{split}
 xb&=p^{-1}bx,\qquad yb=pby,\qquad yx=xy-\vartheta b^2,\\
 a'b&=ba'+t_0(x^2-y^2),\\
 a'x&=p^{-1}xa'+s_0by,\qquad
 a'y=pya'-s_0bx.
\end{split}
\]
検証器\texttt{verify\_cubic\_as4\_ds211\_v.py}は四つのDiamondを全 $p$ で確認し、
PBW基底 $b^i x^j y^k(a')^l$ を得る。

中央多項式変数 $z_1,z_2$ を持つ量子トーラス $xb=p^{-1}bx$ 上で
\[
 y=z_2x^{-1}-p b^2x^{-1},\qquad
 a'=b^{-1}(z_1-\alpha x^2-\beta y^2)
\]
と定める。全六関係に加え、元の $\mathcal V_h$ の全六関係も検証器で確認した。
対称射影の列階数六と新しいPBW二次元数十から、両表示の二次関係空間は一致する。
$z_1,z_2$ の補助重みをそれぞれ三、一とすると、PBW単項式の先頭項は
非零scalar倍の $z_1^l z_2^k b^{i-l}x^{j-k}$ となり、互いに異なる。
従って模型は忠実であり
\[
 z_2=xy+p^{-1}b^2,\qquad z_1=ba'+\alpha x^2+\beta y^2
\]
は実際の中心二次元である。

前節と同じ先頭項の指数比較から、$p^3\ne1$ では中心三次は零である。
また
\[
 h-\frac1{48}=-\frac{p^2+p+1}{12(p-1)^2}
\]
なので、残余は正確に $h=1/48$ である。
$p^3=1$ では量子トーラスのdivision algebraは中心上九次元となる。
$\ell=u a'+v b+w x+t y$ のreduced traceは
\[
 \operatorname{Trd}\ell=0,\qquad
 \tfrac12\operatorname{Trd}(\ell^2)
 =3\bigl(uvz_1+(wt+u^2/4)z_2\bigr).
\]
後式と、非中央Laurent項のtraceが消えることも三乗根の係数環で厳密検算した。
従って
\[
 N_{a'}=(a')^3-\tfrac34z_2a'
\]
は中心三次である。
$b^3,x^3,y^3$ も中心であり、先頭項比較によりこの四本は独立で全中心三次を張る。
実際、中心三次の先頭指数 $(i-l,j-k)$ は三の倍数なので純三乗以外に可能性がなく、
各生成元を引けば先頭項が下がる。

Cayley--Hamiltonは全一次元に対して
\[
 \operatorname{Nrd}(\ell)
 =\ell\Bigl(\ell^2-3\bigl(uvz_1+(wt+u^2/4)z_2\bigr)\Bigr)\in Z(B)_3
\]
を与える。忠実なdivision模型で非零一次元のnormは非零なので、
再び基点のない射影norm $\mathbb P^3\to\mathbb P^3$ の全射性を得る。
全非零中心三次が分解する。
最後に、全 $h$ のAS4実現では元の正規化関係行列と $A_1$ の反巡回最高tensorから
$\mu=I$ であることを独立に確認してある。
従って全H3候補を中心三次へ制限でき、例外値も非整域となる。
\end{proof}
